\exo \textit{[calculer, Raisonner]}

Dans un repère, on donne les vecteurs :
$\overrightarrow{AB}(1;3)$ ; $\overrightarrow{CD}(-2;4)$ et $\overrightarrow{EF}(0;10)$.

\begin{enumerate}
\item Déterminer les coordonnées des vecteurs $\vec{u}=\overrightarrow{AB}+\overrightarrow{CD}$ et $\vec{v}=\overrightarrow{EF}+\overrightarrow{BA}$.
\item Qu'en déduisez-vous?
\end{enumerate}



\exo \textit{[Représenter]}

\begin{enumerate}
\item Reproduire la figure de l'exercice LS 42 p 208 puis représenter une base de vecteurs ($\vec{i}$ ; $\vec{j}$) bien choisie. 
\item Lire les coordonnées du vecteur $\vec{u}$ dans cette base ($\vec{i}$ ; $\vec{j}$).
\item Calculer les coordonnées des vecteurs  $\dfrac{1}{4}\vec{u}$ et $-\dfrac{3}{2}\vec{u}$ puis répondre à la question du manuel.

\end{enumerate}



\exo \textit{[Représenter]}

On considère les vecteurs ci-dessous dans la base ($\vec{i}$ ; $\vec{j}$).

\begin{minipage}{0.6\linewidth}
Dans chaque cas, donner les coordonnées des vecteurs dans la base ($\vec{i}$ ; $\vec{j}$) et en déduire la relation de colinéarité.

\begin{enumerate}[label=\alph*)]
\item $\vec{m}(....;....)$ et $\vec{n}(....;....)$ donc $\vec{n}=....\vec{m}$.
\item $\vec{m}(....;....)$ et $\vec{p}(....;....)$ donc $\vec{m}=....\vec{p}$.
\item $\vec{j}(....;....)$ et $\vec{b}(....;....)$ donc $\vec{b}=....\vec{j}$.
\item $\vec{i}(....;....)$ et $\vec{a}(....;....)$ donc $\vec{i}=....\vec{a}$.
\end{enumerate}
\end{minipage}
\hfill
\begin{minipage}{0.4\linewidth}


\medskip

\begin{center}
\newrgbcolor{wqwqwq}{0.3764705882352941 0.3764705882352941 0.3764705882352941}
\psset{xunit=0.7cm,yunit=0.7cm,algebraic=true,dimen=middle,dotstyle=o,dotsize=5pt 0,linewidth=2pt,arrowsize=3pt 2,arrowinset=0.25}
\begin{pspicture*}(-2.5,-1.5)(5.5,4.5)
\multips(0,-1)(0,1){6}{\psline[linestyle=dashed,linecap=1,dash=1.5pt 1.5pt,linewidth=0.4pt,linecolor=wqwqwq]{c-c}(-2.5,0)(5.5,0)}
\multips(-2,0)(1,0){8}{\psline[linestyle=dashed,linecap=1,dash=1.5pt 1.5pt,linewidth=0.4pt,linecolor=wqwqwq]{c-c}(0,-1.5)(0,4.5)}
\psline[linewidth=2pt]{->}(0,0)(1,0)
\psline[linewidth=2pt]{->}(0,0)(0,1)
\rput[tl](0.24,-0.18){$\vec{i}$}
\rput[tl](-0.54,0.5){$\vec{j}$}
\psline[linewidth=2pt]{->}(1,1)(5,1)
\psline[linewidth=2pt]{->}(-2,3)(-2,0)
\psline[linewidth=2pt]{->}(-1,2)(0,3)
\psline[linewidth=2pt]{->}(1,2)(3,4)
\psline[linewidth=2pt]{->}(4,2)(1,-1)
\psline[linewidth=2pt]{->}(1,2)(3,4)
\rput[tl](-2.4,2.1){$\vec{b}$}
\rput[tl](-0.9,3.1){$\vec{m}$}
\rput[tl](1.82,3.7){$\vec{n}$}
\rput[tl](2.46,0.25){$\vec{p}$}
\rput[tl](2.34,1.7){$\vec{a}$}
\end{pspicture*}
\end{center}
\end{minipage}


\exo \textit{[Calculer]}

\begin{enumerate}
\item Que fait la fonction suivante?

\begin{center}
\begin{minipage}{0.64\linewidth}
\begin{pythoncodenumligne}
def Vect(x_A, y_A, x_B, y_B, x_C, y_C):
    x_AB = x_B - x_A
    y_AB = y_B - y_A
    x_AC = x_C - x_A
    y_AC = y_C - y_A
    return (x_AB, y_AB, x_AC, y_AC)
\end{pythoncodenumligne}
\end{minipage}
\end{center}


\item Modifier la fonction précédente pour qu'elle renvoie VRAI lorsque les points A, B et C sont alignés et FAUX dans le cas contraire.
\item Tester votre fonction en utilisant les points de l'exercice LS 64 p 211 et CE 4 p 48.
\end{enumerate}



\exo \textit{[Calculer, Raisonner]}

\begin{enumerate}
\item Faire l'exercice CE 4 p 45.
\item Placer le point M sur la figure et le point D$(8~;~2)$.
\item Montrer que les points C, D et M sont alignés.
\end{enumerate}


\medskip

\exo \textit{[Calculer, Raisonner]}

Dans un repère, on donne les points : A$(3~;~4)$ ; B$(2~;~1)$ et C$(6~;~1)$.

D est le point tel que $\overrightarrow{AD}=2\overrightarrow{AB}-\overrightarrow{BC}$ et E est le symétrique de C par rapport à B.

\begin{enumerate}
\item Calculer les coordonnées des vecteurs $\overrightarrow{AB}$ et $\overrightarrow{BC}$.
\item En déduire les coordonnées du point D.
\item Faire la figure. Que pouvez-vous conjecturer pour les droites (AB) et (DE)?
\item Prouver votre conjecture.
\end{enumerate}








